\subsection{仿射线性簇}

定义 2. 给定仿射空间 E，E 的子集 V 称为仿射线性簇（或简称线性簇，或 E 的仿射子集），如果对 V 的点的每一族 \((x_{t})_{t\in I}\) 以及 K 中具有有限支集的每一族元素 \((\lambda_t)_{t\in I}\) ，只要 \(\sum_{t\in I}\lambda_t = 1\) ，重心 \(\sum_{t\in I}\lambda_t x_t\) 就属于 V。

这等同于说，定义2的条件对V的每一有限点族都成立。

空集是一个线性簇；线性簇的任意交集仍是一个线性簇。

设 V 是 E 的一个非空子集， \(a\) 是 V 的一个点；关系

\[
x - a = \sum_ {i = 1} ^ {n} \lambda_ {i} (x _ {i} - a)
\]

意味着 \(x\) 是族的一个重心 \(\sum_{i=1}^{n} \lambda_i x_i + \left(1 - \sum_{i=1}^{n} \lambda_i\right)a\) ，该族由诸 \(x_i\) 与 \(a\) 组成。因此：

命题 2. 对于仿射空间 E 的非空子集 V，要使 V 成为线性簇，当且仅当在取 V 中一点为原点所得到的 E 的向量空间结构下，V 是向量子空间。

特别地，向量空间 T（视为仿射空间）的非空仿射线性簇恰好是 T 的向量子空间的平移；因此，T 的向量子空间就是包含 0 的线性簇。

设 V 是仿射空间 E 的非空线性簇；当 x 和 y 遍历 V 时，自由向量 x - y 的集合是 E 的平移空间 T 的一个向量子空间 D，称为 V 的方向：因为，若 \(a \in V\) ，则

\[
x - y = (x - a) - (y - a)
\]

于是我们的断言由命题2得出。显然，D在V上忠实且传递地作用，因此V典范地具有与D相联系的仿射空间结构。所谓仿射簇V的维数，我们指的是V在此仿射空间结构下的维数，即向量空间D的维数。维数为0的线性簇是E的点；维数为1（相应地，2）的线性簇称为E的直线（相应地，平面）。

每个向量 \(\neq0\) 属于某直线的方向，称为该直线的一个方向向量；其关于T的一组基的分量构成所谓该直线的方向参数组。

线性簇V在E中的余维数是指其方向D在T中的余维数；E中余维数为1的线性簇称为E的一个（仿射）超平面。

具有相同方向的两个线性簇称为平行的；这等同于说其中一个是由另一个通过平移得到的。若V是T（视为仿射空间）中的线性簇，则其方向是与V平行且包含0的线性簇。

命题 3. 给定一族点 \((a_{t})_{t\in I}\) （属于一个仿射空间 \(\mathbf{E}\) ），集合 \(\mathbf{V}\) ——由重心 \(\sum_{i\in I}\lambda_i a_i\) (( \(\lambda_i\) ) 具有有限支集， \(\sum_{i\in I}\lambda_i = 1\) ) 组成—— 是 \(\mathbf{E}\) 的一个线性簇。

若族 \((a_{t})\) 为空，则 \(\mathbf{V} = \varnothing\) ，这是由条件 \(\sum_{i}\lambda_{i} = 1\) 所致。因此可设族 \((a_{t})\) 非空，此时在 E 中取某个 \(a_{t}\) 作为原点，命题便是显然的。

线性簇 V 显然是包含诸 \(a_{t}\) 的最小线性簇；称它为由族 \((a_{t})\) 生成的线性簇，并称该族为 V 的一个生成系。

在命题 3 的记号下，假设族 \((a_{\iota})\) 非空，要使每个点 \(x \in V\) 的形如 \(x = \sum_{\iota} \lambda_{\iota} a_{\iota}\) 的表达式是唯一的，必要且充分条件是：用 \(\kappa\) 表示 I 的任意指标时，向量族 \(a_{\iota} - a_{\kappa}\) （其中 \(\iota\) 遍历指标集 \(\neq \kappa\) ）在 T 中是自由的。这时称 E 的点的族 \((a_{\iota})_{\iota \in I}\) 是仿射自由的（或称其元素构成一个仿射自由系，或称它们仿射无关），并称 \(\lambda_{\iota}\) 是 \(x\) 的指标为 \(\iota\) 的、关于仿射自由族 \((a_{\iota})\) 的重心坐标。

设 \((a_{i})_{i\in I}\) 是 E 中一族点，若它不是仿射自由的，则称它是仿射相关的。

命题 4. 要使点的非空族 \((a_{t})_{t\in I}\) 在仿射空间 \(E\) 中是仿射相关的，必要且充分条件是：存在一个族 \((\lambda_t)_{t\in I}\) ，其元素在 \(K\) 中且不全为零，具有有限支集，使得 \(\sum_{i\in I}\lambda_i = 0\) 且 \(\sum_{i\in I}\lambda_i a_i = 0\) 。

给定指标 \(\kappa \in I\) ，说向量族 \((a_{t} - a_{\kappa})\) （其中 \(t\) 遍历指标集 \(\neq \kappa\) ）在 \(T\) 中相关，意味着存在一个不全为零的标量族 \((\lambda_{t})_{t \neq \kappa}\) ，使得 \(\sum_{t \neq \kappa} \lambda_{t}(a_{t} - a_{\kappa}) = 0\) ，它也可写成 \(\sum_{t \in I} \lambda_{t} a_{t} = 0\) ，其中 \(\lambda_{\kappa} = -\sum_{t \neq \kappa} \lambda_{t}\) ，换言之 \(\sum_{t \in I} \lambda_{t} = 0\) 。

命题 5. 要使点的非空族 \((a_{t})_{t\in I}\) 在仿射空间 \(E\) 中是仿射自由的，必要且充分条件是：对每个指标 \(\kappa\in I\) ， \(a_{\kappa}\) 不属于由诸 \(a_{t}\) （指标 \(\neq\kappa\) ）生成的线性簇。

当 I 只有一个元素时，命题显然成立。否则，在 E 中取某个 \(a_{t}\) （指标 \(\neq\kappa\) ）作为原点，命题即由 §7 第 1 号的注记得出。

